\section{Related Work}
\label{sec:related-work}


\paragraph{Lens models.} The thin lens
\cite{Hecht+2023,pharr2023physically} gives the perspective camera a
finite aperture, so that a scene point off the plane of focus is
imaged to a circle of confusion. It is tied to the perspective projection and has
no natural counterpart in classic distortion models \cite{kannala-brandt, double-sphere, Brown1966DecenteringDO, Conrady1919DecentredL}. The plane of focus and the vignetting that emerge from this model are too simple to capture the actual phenomena. The
surface of best focus of a real lens is curved, which is known in
optics as field curvature \cite{Kingslake2010}. Tang and Kutulakos
\cite{tang2013aberrated} derive the blur of a rotationally symmetric
lens as an integral over its aperture and show that field curvature
acts as a defocus that depends on the position in the image. This is
what a focus surface describes, and it justifies our design, which
replaces their fixed, symmetric dependence by a free-form one. Their
analysis is theoretical, and the aberrations for which the rays of a
sensor position no longer meet in a point are not part of our model. Models that trace the
rays through the lens elements \cite{kolb1995realistic} need the
optical design of the lens, which manufacturers rarely publish. The
neural lens model of Xian \etal \cite{xian2023neurallens} replaces the
distortion of a pinhole camera by an invertible network and models the
vignetting as a gain per pixel of the image, without an aperture. The thin lens is also the only finite aperture of the
differentiable renderer Mitsuba
\cite{NimierDavidVicini2019Mitsuba2,Mitsuba3}. Its cameras are
perspective and without lens distortion, and the aperture radius and
the focus distance of its thin lens are not differentiable parameters,
so neither the projection of a general central model nor its aperture
and focus can be calibrated with it. Our calibration keeps every
camera parameter differentiable: neither a parameter-dependent
discontinuity nor a parameter-dependent integration domain remains in
its formulation.

\paragraph{Radiometric calibration and image corrections.} Radiometric calibration methods
estimate the vignetting, often together with the response function,
from the pixel values of overlapping images
\cite{Kim-PAMI-2008,single-image-vignetting,direct-sparse-odometry}.
In \cite{Kim-PAMI-2008,xian2023neurallens} the vignetting is an empirical function of the image position, and
each pixel value is related to the radiance of a single scene point.
Reconstruction methods likewise add corrections around a geometric
camera, such as per-image gains \cite{brown2007automatic}, appearance
embeddings \cite{martinbrualla2021nerfw} or blur kernels
\cite{ma2022deblurnerf}. These compensate for the unmodeled radiometry in the image, where
defocus, which depends on the depth of the scene point, cannot be
corrected consistently. The radiometric camera models it inside the
camera, as an integral over its aperture.

\paragraph{Calibration from pixel values.} Kang and Weiss
\cite{kang2000calibrate} obtain a rough focal length of a perspective
camera from the $\cos^4$ falloff in a single image of a uniform
surface. Rehder \etal \cite{rehder2017direct} formulate a
maximum-likelihood calibration on the image intensities by rendering
the target. They assume an evenly lit target and no vignetting, and
take a single radiance sample per pixel, to which they attribute that
their static calibration is no more precise than one from features;
defocus, vignetting and uneven lighting are named as future work.
Hannemose \etal \cite{Hannemose_2019} also render the target densely,
with a pinhole camera with Brown--Conrady distortion, but their
calibration is photometric and not radiometric: the rendering blurs
the texture with a Gaussian whose width is a free parameter of each
interest point, and it models neither the lighting nor the intensity
scale of the image. A separate line calibrates the
blur of a lens as a learned point spread function, apart from the
projection of the camera \cite{lin2023blurfields}. In the radiometric
camera, the defocus follows from the aperture and the focus surface of
the same model that projects the scene, for any central projection,
and it is calibrated together with the vignetting, the lighting and
the target.

% Ties's TODO, addressed by the paragraphs above (11-17):
% Discuss related work w.r.t. our contributions ! Why cant Mitsuba with thin lens just do what we are proposing? Where is our formulation different from PBR-book and why. Or we just implemented something that has been known for a long time ?
% As compared to Hannemose_2019: what is it they cant do ?
% Does it really help to have unsharpness/vignetting (this must be shown in the experiments then)

% \textbf{I (Ties) would say it is as follows:\\
% Camera sensors measure radiometric properties, not geometric properties\\
% $\Rightarrow$ A bias in the camera model propagates through the subsequent computations. a) the reduction of the sensor measurements to geometric measurements by applying a detector leads to a higher reprojection error in calibration, b) without modeling the aperture and defocus, the defocus will be absorbed \eg in texture learning\\
% $\Rightarrow$ There exists \eg the thin lens model, the thick lens model and very accurate models which model every lens intersection\\
% $\qquad\Rightarrow$ Thin lens and thick lens don't model distortion; the complex models are too complex\\
% $\Rightarrow$ we need a generic lens model which includes distortion, vignetting and defocus\\
% $\Rightarrow$ Using thin lens + distortion is not sufficient because the focus surface of a thin lens is per definition a plane, but for a fisheye for example it is more like a sphere, which is shown in \cref{fig:focus-surfaces}\\
% In total, we need:\\
% I: Distortion $\Rightarrow$ We can use every distortion model in our pipeline (Geometric accuracy experiment)\\
% II: Aperture for defocus $\Rightarrow$ A camera model which does not model blur introduces a bias in inverse rendering (Texture learning experiment)\\
% III: A focus surface $\Rightarrow$ The focus surfaces differ between lenses (\cref{fig:focus-surfaces})\\
% IV: Vignetting $\Rightarrow$ Not modeling the vignetting will also introduce a bias in inverse rendering (we do not show that) and is useful for example in panorama stitching (Vignetting experiment)
% }





% Moved into the introduction (the paragraph on radiometric extensions of the camera), 11-17:
% \subsection{Photometric Camera Calibration}
% Hannemose et al. (2019) \cite{Hannemose_2019} perform a photometric camera calibration. Therefore, they blur the texture of the calibration board, project the board into a 2D image and minimize the difference between the pixel values of the original captured calibration board and the rendered one. This method does not account for lighting, which is crucial when using a photometric or radiometric loss.
